\subsection{Decreasing approximation of bounded potentials}
\label{subsec:hessian-bounded-approximation}
\index{Admissible potential!decreasing smooth approximation}

Upper semicontinuity gives smooth obstacles decreasing to a bounded
potential. The same penalty equation gives admissible functions between
the potential and each obstacle, up to an error tending to zero.
Choosing successive obstacles below the preceding smooth approximation
makes the admissible approximations decrease as well.

\begin{lemma}\label{hessian:smooth-majorants}
Every bounded upper semicontinuous real function $u$ on $X$ is the
pointwise decreasing limit of smooth functions $h_j>u$.
\end{lemma}
\begin{proof}
Fix a Riemannian distance $d$ on $X$ and put
\[
 c_j(x)=\sup_{y\in X}\{u(y)-j d(x,y)\}.
\]
Each $c_j$ is Lipschitz, $c_j\geq u$, and $c_{j+1}\leq c_j$.
The supremum is attained. If $y_j$ attains it at a fixed $x$, then
$j d(x,y_j)\leq\sup_X u-u(x)$, so $y_j\to x$.
Upper semicontinuity gives $\limsup_j c_j(x)\leq u(x)$, and hence
$c_j\downarrow u$ pointwise.

Set $\delta_j=4^{-j}$. Smooth approximation of continuous functions
on a compact manifold gives $s_j\in C^\infty(X)$ with
$\|s_j-c_j\|_\infty\leq\delta_j/4$.
Take $h_j=s_j+\delta_j$. Then
\[
 u+\tfrac34\delta_j\leq c_j+\tfrac34\delta_j\leq h_j
                     \leq c_j+\tfrac54\delta_j,
 \qquad
 h_{j+1}\leq c_j+\tfrac5{16}\delta_j
                    < c_j+\tfrac34\delta_j\leq h_j.
\]
The bounds also imply $h_j(x)\to u(x)$ at every point.
\end{proof}

\Needspace{16\baselineskip}
\begin{lemma}\label{hessian:bounded-penalty-squeeze}
Let $u$ be bounded admissible and $h\geq u$ smooth. Choose
\[
 b>\max\{1,\|H_m(h)/\omega^n\|_\infty\}.
\]
For each positive integer $N$, let $w_N$ be the smooth strictly
admissible solution of
\[
 H_m(w_N)=b e^{N(w_N-h)}\omega^n.
\]
With
\[
 M_N=\|h\|_\infty+\frac{\log b}{N},\qquad
 E_N=\frac{m\log(N+1)+\log b}{N},
\]
one has $\|w_N\|_\infty\leq M_N$ and
\[
 u-E_N-2M_N/N\leq w_N\leq h.
\]
\end{lemma}
\begin{proof}
Existence follows from \cref{hessian:smooth-proper-existence}. The
bounds are \eqref{eq:hessian-penalty-squeeze} and the zero-order bound in
the proof of \cref{hessian:continuous-smooth-approximation}, which use
only that $u$ is bounded admissible and $u\leq h$.
\end{proof}

\begin{theorem}\label{hessian:bounded-decreasing-approximation}
Every bounded admissible potential $u$ on $X$ is the pointwise
decreasing limit of smooth strictly admissible potentials.
The sequence may be chosen uniformly bounded.
\end{theorem}
\begin{proof}
Choose smooth $h_k\downarrow u$ by \cref{hessian:smooth-majorants}
and put $\varepsilon_j=8^{-j}$. We construct recursively smooth strictly
admissible $u_j$ with $u_j\geq u+\varepsilon_j$.
For $j=1$, take $k_1=1$. At a subsequent step choose
$k_j\geq\max\{j,k_{j-1}+1\}$ so large that
\[
 h_{k_j}<u_{j-1}-\tfrac12\varepsilon_{j-1}\quad\hbox{on }X.
\]
Such an index exists: the open sets
$\{h_k<u_{j-1}-\varepsilon_{j-1}/2\}$ increase with $k$ and cover
$X$, since $h_k\downarrow u\leq u_{j-1}-\varepsilon_{j-1}$.
Compactness gives a single index for all of $X$.

For this fixed obstacle $h_{k_j}$ choose $b_j$ as in
\cref{hessian:bounded-penalty-squeeze}, and then choose $N_j$ so large
that its error $E_{N_j}+2M_{N_j}/N_j$ is at most $\varepsilon_j$.
The resulting smooth solution $w_j$ satisfies
\[
 u-\varepsilon_j\leq w_j\leq h_{k_j}.
\]
Put $u_j=w_j+2\varepsilon_j$. Adding a constant preserves strict
admissibility, and $u_j\geq u+\varepsilon_j$. For $j\geq2$,
\[
 u_j\leq h_{k_j}+2\varepsilon_j
 <u_{j-1}-\tfrac12\varepsilon_{j-1}+2\varepsilon_j
 =u_{j-1}-\tfrac14\varepsilon_{j-1}.
\]
Finally, $u\leq u_j\leq h_{k_j}+2\varepsilon_j$ and $k_j\geq j$
give pointwise convergence to $u$. The bounds
$\inf_X u\leq u_j\leq\sup_X u_1$ give a common uniform bound.
\end{proof}
